
\section{Optimización y Experimentación}

\subsection{Consultas SQL para el experimento}

\subsubsection{Descripción de las 5 consultas intermedias seleccionadas (no triviales)}

Las cinco consultas trabajan sobre \texttt{LecturaVibracion}, la tabla que crece
con la escala, y responden preguntas de mantenimiento predictivo.

\textbf{Consulta 1:} ¿Qué lecturas anómalas (RMS $\ge 4$) superan en 1.5 veces o
más el promedio de las 24 horas previas de su mismo sensor?

\textit{Justificación:} detecta picos que no son ruido aislado sino un salto
frente al historial reciente del sensor, y los ubica en su turbina y parque.

\textbf{Consulta 2:} En las últimas 48 horas, ¿qué sensores vibran más que el
promedio de su parque? (top 5 por parque)

\textit{Justificación:} compara cada sensor contra sus pares del mismo parque
para priorizar la revisión de los que se salen del comportamiento conjunto.

\textbf{Consulta 3:} ¿Qué turbinas tienen más rachas de tres o más lecturas
consecutivas con RMS $> 2.5$? (top 10)

\textit{Justificación:} una lectura alta aislada puede ser ruido; una racha
sostenida indica deterioro real del componente.

\textbf{Consulta 4:} ¿Cómo venía la vibración en las cinco lecturas previas a
cada alerta grave? (por parque y componente)

\textit{Justificación:} mide la tendencia (pendiente) antes de la alerta para
distinguir fallas progresivas de eventos súbitos.

\textbf{Consulta 5:} ¿Qué turbinas pasan más rápido de ``normal'' a ``alerta''?
(top 10, incluyendo alertas sin orden de trabajo)

\textit{Justificación:} cuantifica la velocidad de degradación y detecta
alertas que todavía no tienen una orden asociada.

\subsubsection{Implementación de consultas en SQL}

\textbf{Consulta 1: Lecturas anómalas frente a su historial de 24 h}
\lstinputlisting[language=SQL]{consultas_limpias/c1.sql}

\textbf{Consulta 2: Sensores por encima del promedio de su parque (48 h)}
\lstinputlisting[language=SQL]{consultas_limpias/c2.sql}

\textbf{Consulta 3: Rachas de vibración alta por turbina}
\lstinputlisting[language=SQL]{consultas_limpias/c3.sql}

\textbf{Consulta 4: Tendencia previa a las alertas graves}
\lstinputlisting[language=SQL]{consultas_limpias/c4.sql}

\textbf{Consulta 5: Velocidad de degradación de cada turbina}
\lstinputlisting[language=SQL]{consultas_limpias/c5.sql}

\subsection{Metodología del experimento}

Para que la comparación sea válida, la variante \emph{sin índices} no se limita
a borrar los índices creados por el proyecto: también se desactiva el uso de los
índices automáticos de llave primaria y \texttt{UNIQUE} en la sesión, porque
PostgreSQL los aprovechaba y la prueba dejaba de ser ``sin índices''. El
procedimiento fue:

\begin{enumerate}
    \item \textbf{Preparación:} se parte de las bases \texttt{bd\_1k},
    \texttt{bd\_10k}, \texttt{bd\_100k} y \texttt{bd\_1m}, generadas con
    \texttt{crear\_todo.sh}, y se ejecuta \texttt{ANALYZE} para actualizar las
    estadísticas del planificador.

    \item \textbf{Sin índices:} se eliminan todos los índices \texttt{idx\_*} y,
    en la misma sesión, se ejecuta
    \begin{lstlisting}[language=SQL]
SET enable_indexscan = off;
SET enable_indexonlyscan = off;
SET enable_bitmapscan = off;
    \end{lstlisting}
    Así el planificador solo puede usar \texttt{Seq Scan}. El programa verifica
    en cada plan que la variante ``sin'' no use ningún índice.

    \item \textbf{Calentamiento:} antes de medir se ejecuta la consulta una vez
    para llevar los datos a la caché compartida.

    \item \textbf{Repeticiones:} se ejecutan 5 veces por consulta, escala y
    variante. Cada valor de las tablas es una ejecución; el promedio y la
    desviación estándar se calculan sobre esas 5 corridas. Si una corrida supera
    los 120\,s se mide una sola vez ($n=1$) para no extender el experimento; es
    el caso de C1 y C5 con 1M sin índices.

    \item \textbf{Con índices:} se ejecuta \texttt{indices/crear\_indices.sql}
    (índices elegidos por teoría) y se repite la misma medición en una sesión
    limpia.

    \item \textbf{Medición:} todo se hace con \texttt{benchmark.py} en una sola
    sesión de \texttt{psycopg} por variante, registrando el tiempo de ejecución
    y el tipo de acceso del plan. Los planes que se muestran en el informe se
    obtienen con \texttt{EXPLAIN (ANALYZE, BUFFERS)} en pgAdmin.
\end{enumerate}

\subsection{Optimización de consultas}

La regla usada para elegir cada índice es la siguiente: la igualdad
($=$, \texttt{IN}, join por llave) va con \textbf{Hash}; el rango, la
desigualdad y el orden ($<$, $>$, $\ge$, \texttt{ORDER BY},
\texttt{PARTITION BY \dots ORDER BY}) van con \textbf{B+ tree}; si hay igualdad
y rango en varias columnas se usa B+ tree, porque Hash solo indexa una columna y
solo sirve para igualdad. En PostgreSQL, \texttt{CREATE INDEX} sin
\texttt{USING} crea un B-tree, que es un B+ tree.

\subsubsection{Planes de índices para Consulta 1}
\begin{lstlisting}[language=SQL]
CREATE INDEX idx_btree_lectura_sensor_fecha
  ON LecturaVibracion (id_sensor, fecha_hora) INCLUDE (valor_rms);

CREATE INDEX idx_btree_lectura_rms
  ON LecturaVibracion (valor_rms);
\end{lstlisting}
La subconsulta correlacionada busca \texttt{(id\_sensor, fecha\_hora)} en un
rango de 24 h (B+ tree) y el filtro \texttt{valor\_rms >= 4} usa el B+ tree
sobre \texttt{valor\_rms}.

\subsubsection{Planes de índices para Consulta 2}
\begin{lstlisting}[language=SQL]
CREATE INDEX idx_btree_lectura_fecha
  ON LecturaVibracion (fecha_hora) INCLUDE (id_sensor, valor_rms, valor_pico);
\end{lstlisting}
El predicado es un rango de 48 h sobre \texttt{fecha\_hora}; el \texttt{INCLUDE}
permite un \texttt{Index Only Scan} sin volver a la tabla.

\subsubsection{Planes de índices para Consulta 3}
\begin{lstlisting}[language=SQL]
CREATE INDEX idx_btree_lectura_sensor_fecha
  ON LecturaVibracion (id_sensor, fecha_hora) INCLUDE (valor_rms);
\end{lstlisting}
Las ventanas necesitan las lecturas ordenadas por sensor y fecha; el B+ tree las
entrega en ese orden (evita el \texttt{Sort} de toda la tabla).

\subsubsection{Planes de índices para Consulta 4}
\begin{lstlisting}[language=SQL]
CREATE INDEX idx_btree_lectura_sensor_fecha
  ON LecturaVibracion (id_sensor, fecha_hora) INCLUDE (valor_rms);

CREATE INDEX idx_hash_alerta_nivel
  ON AlertaMantenimiento USING hash (nivel);
\end{lstlisting}
El \texttt{JOIN LATERAL} busca las 5 lecturas anteriores de un sensor
(igualdad + orden) y el filtro \texttt{nivel IN ('Alta','Critica')} es igualdad
(Hash).

\subsubsection{Planes de índices para Consulta 5}
\begin{lstlisting}[language=SQL]
CREATE INDEX idx_btree_lectura_sensor_fecha
  ON LecturaVibracion (id_sensor, fecha_hora) INCLUDE (valor_rms);

CREATE INDEX idx_hash_alertaorden_alerta
  ON AlertaOrden USING hash (id_alerta);
\end{lstlisting}
Cada alerta busca la última lectura normal anterior a una fecha (igualdad +
rango) y el \texttt{NOT EXISTS} resuelve \texttt{ao.id\_alerta = a.id\_alerta}
por igualdad.

\paragraph{Índices de las llaves foráneas.} Los joins por llave foránea se apoyan en índices Hash:
\texttt{SensorVibracion(id\_componente)},
\texttt{ComponenteTurbina(id\_turbina)} y \texttt{Turbina(id\_parque)}. El
conjunto completo de índices está en el repositorio del proyecto.

\subsection{Plataforma de Pruebas}

\begin{table}[H]
\centering
\small
\caption{Hardware y software utilizado}
\begin{tabular}{|l|l|}
\hline
\textbf{Componente} & \textbf{Descripción} \\ \hline
Memoria física instalada (RAM) & 12 GB \\ \hline
Procesador & Intel Core i5-1035G4 @ 1.10 GHz (4 núcleos / 8 hilos) \\ \hline
Almacenamiento & SSD NVMe SK Hynix 512 GB \\ \hline
Sistema Operativo & Fedora Linux 44 Workstation (kernel 7.3.3) \\ \hline
pgAdmin & pgAdmin 4 v9.19 \\ \hline
PostgreSQL & PostgreSQL 16.15 \\ \hline
\end{tabular}
\end{table}

\subsection{Medición de tiempos}

Se registra el \textit{Execution Time} de \texttt{EXPLAIN ANALYZE} en
milisegundos. Cada tabla muestra las cinco ejecuciones por escala, su promedio
y su desviación estándar.

\subsubsection{Sin índices}

\begin{table}[H]
\centering
\small
\caption{Tiempos de ejecución --- Consulta 1 --- Sin índices}
\begin{tabular}{|c|r|r|r|r|}
\hline
\textbf{Ejecución} & \textbf{1k (ms)} & \textbf{10k (ms)} & \textbf{100k (ms)} & \textbf{1M (ms)} \\ \hline
1 & 0.375 & 23.43 & 2\,178.3 & 197\,558.0 \\ \hline
2 & 0.344 & 23.50 & 2\,175.0 & \textit{--} \\ \hline
3 & 0.328 & 24.49 & 2\,033.0 & \textit{--} \\ \hline
4 & 0.311 & 25.32 & 1\,882.1 & \textit{--} \\ \hline
5 & 0.366 & 24.18 & 1\,871.4 & \textit{--} \\ \hline
\textbf{Promedio} & 0.345 & 24.18 & 2\,028.0 & 197\,558.0 \\ \hline
\textbf{Desviación estándar} & 0.026 & 0.775 & 150.03 & 0.000 \\ \hline
\end{tabular}
\par\footnotesize En 1M la corrida sin índices superó el umbral de 120\,s, por lo que se midió una sola vez ($n=1$); las celdas restantes no aplican.
\end{table}

\begin{table}[H]
\centering
\small
\caption{Tiempos de ejecución --- Consulta 2 --- Sin índices}
\begin{tabular}{|c|r|r|r|r|}
\hline
\textbf{Ejecución} & \textbf{1k (ms)} & \textbf{10k (ms)} & \textbf{100k (ms)} & \textbf{1M (ms)} \\ \hline
1 & 3.526 & 14.87 & 59.23 & 167.53 \\ \hline
2 & 3.574 & 15.22 & 59.17 & 167.03 \\ \hline
3 & 3.375 & 13.84 & 59.02 & 161.30 \\ \hline
4 & 3.302 & 13.75 & 59.49 & 167.33 \\ \hline
5 & 3.705 & 14.10 & 62.93 & 166.65 \\ \hline
\textbf{Promedio} & 3.496 & 14.36 & 59.97 & 165.97 \\ \hline
\textbf{Desviación estándar} & 0.160 & 0.656 & 1.665 & 2.629 \\ \hline
\end{tabular}

\end{table}

\begin{table}[H]
\centering
\small
\caption{Tiempos de ejecución --- Consulta 3 --- Sin índices}
\begin{tabular}{|c|r|r|r|r|}
\hline
\textbf{Ejecución} & \textbf{1k (ms)} & \textbf{10k (ms)} & \textbf{100k (ms)} & \textbf{1M (ms)} \\ \hline
1 & 1.390 & 13.04 & 147.75 & 1\,639.1 \\ \hline
2 & 1.384 & 12.86 & 147.31 & 1\,671.0 \\ \hline
3 & 1.391 & 13.69 & 149.32 & 1\,643.4 \\ \hline
4 & 1.471 & 13.21 & 149.03 & 1\,638.2 \\ \hline
5 & 1.477 & 12.97 & 148.74 & 1\,637.9 \\ \hline
\textbf{Promedio} & 1.423 & 13.15 & 148.43 & 1\,645.9 \\ \hline
\textbf{Desviación estándar} & 0.047 & 0.324 & 0.863 & 14.19 \\ \hline
\end{tabular}

\end{table}

\begin{table}[H]
\centering
\small
\caption{Tiempos de ejecución --- Consulta 4 --- Sin índices}
\begin{tabular}{|c|r|r|r|r|}
\hline
\textbf{Ejecución} & \textbf{1k (ms)} & \textbf{10k (ms)} & \textbf{100k (ms)} & \textbf{1M (ms)} \\ \hline
1 & 0.445 & 10.46 & 1\,102.7 & 90\,218.5 \\ \hline
2 & 0.340 & 10.71 & 1\,101.4 & 89\,129.8 \\ \hline
3 & 0.425 & 10.46 & 1\,150.9 & 95\,441.7 \\ \hline
4 & 0.331 & 11.43 & 1\,107.8 & 94\,191.3 \\ \hline
5 & 0.339 & 10.46 & 1\,111.3 & 93\,882.1 \\ \hline
\textbf{Promedio} & 0.376 & 10.70 & 1\,114.8 & 92\,572.7 \\ \hline
\textbf{Desviación estándar} & 0.054 & 0.419 & 20.53 & 2\,736.9 \\ \hline
\end{tabular}

\end{table}

\begin{table}[H]
\centering
\small
\caption{Tiempos de ejecución --- Consulta 5 --- Sin índices}
\begin{tabular}{|c|r|r|r|r|}
\hline
\textbf{Ejecución} & \textbf{1k (ms)} & \textbf{10k (ms)} & \textbf{100k (ms)} & \textbf{1M (ms)} \\ \hline
1 & 0.469 & 74.80 & 5\,955.2 & 575\,110.5 \\ \hline
2 & 0.460 & 74.50 & 5\,910.3 & \textit{--} \\ \hline
3 & 0.455 & 75.06 & 5\,903.5 & \textit{--} \\ \hline
4 & 0.455 & 74.84 & 5\,906.8 & \textit{--} \\ \hline
5 & 0.451 & 74.90 & 5\,989.5 & \textit{--} \\ \hline
\textbf{Promedio} & 0.458 & 74.82 & 5\,933.1 & 575\,110.5 \\ \hline
\textbf{Desviación estándar} & 0.007 & 0.203 & 37.93 & 0.000 \\ \hline
\end{tabular}
\par\footnotesize En 1M la corrida sin índices superó el umbral de 120\,s, por lo que se midió una sola vez ($n=1$); las celdas restantes no aplican.
\end{table}


\subsubsection{Con índices}

\begin{table}[H]
\centering
\small
\caption{Tiempos de ejecución --- Consulta 1 --- Con índices}
\begin{tabular}{|c|r|r|r|r|}
\hline
\textbf{Ejecución} & \textbf{1k (ms)} & \textbf{10k (ms)} & \textbf{100k (ms)} & \textbf{1M (ms)} \\ \hline
1 & 0.055 & 0.556 & 5.277 & 55.29 \\ \hline
2 & 0.055 & 0.519 & 4.837 & 55.22 \\ \hline
3 & 0.050 & 0.538 & 4.921 & 55.27 \\ \hline
4 & 0.049 & 0.518 & 5.061 & 57.59 \\ \hline
5 & 0.101 & 0.514 & 4.853 & 56.55 \\ \hline
\textbf{Promedio} & 0.062 & 0.529 & 4.990 & 55.98 \\ \hline
\textbf{Desviación estándar} & 0.022 & 0.018 & 0.183 & 1.058 \\ \hline
\end{tabular}

\end{table}

\begin{table}[H]
\centering
\small
\caption{Tiempos de ejecución --- Consulta 2 --- Con índices}
\begin{tabular}{|c|r|r|r|r|}
\hline
\textbf{Ejecución} & \textbf{1k (ms)} & \textbf{10k (ms)} & \textbf{100k (ms)} & \textbf{1M (ms)} \\ \hline
1 & 3.613 & 13.25 & 47.04 & 50.01 \\ \hline
2 & 3.378 & 13.07 & 47.55 & 50.99 \\ \hline
3 & 3.472 & 13.94 & 48.75 & 50.35 \\ \hline
4 & 3.340 & 13.84 & 48.15 & 51.31 \\ \hline
5 & 3.506 & 14.06 & 48.68 & 50.03 \\ \hline
\textbf{Promedio} & 3.462 & 13.63 & 48.03 & 50.54 \\ \hline
\textbf{Desviación estándar} & 0.108 & 0.443 & 0.738 & 0.584 \\ \hline
\end{tabular}

\end{table}

\begin{table}[H]
\centering
\small
\caption{Tiempos de ejecución --- Consulta 3 --- Con índices}
\begin{tabular}{|c|r|r|r|r|}
\hline
\textbf{Ejecución} & \textbf{1k (ms)} & \textbf{10k (ms)} & \textbf{100k (ms)} & \textbf{1M (ms)} \\ \hline
1 & 1.458 & 13.55 & 144.93 & 1\,438.2 \\ \hline
2 & 1.353 & 13.61 & 145.69 & 1\,460.0 \\ \hline
3 & 1.448 & 13.60 & 144.44 & 1\,457.7 \\ \hline
4 & 1.430 & 13.70 & 143.64 & 1\,430.6 \\ \hline
5 & 1.287 & 13.60 & 144.24 & 1\,433.0 \\ \hline
\textbf{Promedio} & 1.395 & 13.61 & 144.59 & 1\,443.9 \\ \hline
\textbf{Desviación estándar} & 0.073 & 0.057 & 0.772 & 13.96 \\ \hline
\end{tabular}

\end{table}

\begin{table}[H]
\centering
\small
\caption{Tiempos de ejecución --- Consulta 4 --- Con índices}
\begin{tabular}{|c|r|r|r|r|}
\hline
\textbf{Ejecución} & \textbf{1k (ms)} & \textbf{10k (ms)} & \textbf{100k (ms)} & \textbf{1M (ms)} \\ \hline
1 & 0.263 & 0.442 & 3.698 & 38.30 \\ \hline
2 & 0.261 & 0.444 & 3.594 & 38.05 \\ \hline
3 & 0.228 & 0.463 & 3.456 & 39.02 \\ \hline
4 & 0.199 & 0.431 & 3.441 & 33.76 \\ \hline
5 & 0.307 & 0.446 & 3.571 & 32.52 \\ \hline
\textbf{Promedio} & 0.252 & 0.445 & 3.552 & 36.33 \\ \hline
\textbf{Desviación estándar} & 0.041 & 0.012 & 0.106 & 2.968 \\ \hline
\end{tabular}

\end{table}

\begin{table}[H]
\centering
\small
\caption{Tiempos de ejecución --- Consulta 5 --- Con índices}
\begin{tabular}{|c|r|r|r|r|}
\hline
\textbf{Ejecución} & \textbf{1k (ms)} & \textbf{10k (ms)} & \textbf{100k (ms)} & \textbf{1M (ms)} \\ \hline
1 & 0.061 & 0.722 & 5.695 & 62.10 \\ \hline
2 & 0.059 & 0.618 & 5.576 & 65.74 \\ \hline
3 & 0.059 & 0.823 & 5.521 & 84.06 \\ \hline
4 & 0.058 & 0.632 & 5.666 & 77.21 \\ \hline
5 & 0.062 & 1.084 & 5.570 & 70.12 \\ \hline
\textbf{Promedio} & 0.060 & 0.776 & 5.606 & 71.84 \\ \hline
\textbf{Desviación estándar} & 0.002 & 0.191 & 0.072 & 8.848 \\ \hline
\end{tabular}

\end{table}


\subsection{Resultados}

Para cada consulta se muestran el árbol del plan (\textit{Query Plan Tree}) y el
plan de ejecución en las cuatro escalas, sin y con índices, seguidos de su
explicación.

\subsubsection{Consulta 1}
\resultadoconsulta{1}

\textbf{Sin índices:} por cada lectura anómala hay que promediar las 24 horas
previas de su sensor, y cada búsqueda recorre la tabla completa. El costo crece
de forma cuadrática: 0.345 ms (1k), 24.18 ms (10k), 2\,028.0 ms (100k) y
197\,558.0 ms (1M, $n=1$).

\textbf{Con índices:} el B+ tree \texttt{(id\_sensor, fecha\_hora)} resuelve
cada ventana de 24 h en un salto y el índice sobre \texttt{valor\_rms} filtra
las lecturas anómalas. El tiempo baja a 0.062, 0.529, 4.990 y 55.98 ms; con
1M es aproximadamente 3\,529 veces más rápida. La mejora es real en todas las
escalas.


\subsubsection{Consulta 2}
\resultadoconsulta{2}

\textbf{Sin índices:} el rango de 48 h se resuelve con un recorrido secuencial
de \texttt{LecturaVibracion}: 3.496 ms (1k), 14.36 ms (10k), 59.97 ms
(100k) y 165.97 ms (1M).

\textbf{Con índices:} el B+ tree sobre \texttt{fecha\_hora} con
\texttt{INCLUDE} permite un \texttt{Index Only Scan} y reduce el tiempo a
3.462, 13.63, 48.03 y 50.54 ms. Con 1k y 10k la mejora está dentro del
ruido (-1.0\,\% y -5.1\,\%), porque 48 h abarcan casi toda la tabla; con 100k y 1M el
rango ya es selectivo y la ganancia crece hasta -69.5\,\% en 1M.


\subsubsection{Consulta 3}
\resultadoconsulta{3}

\textbf{Sin índices:} la consulta no tiene filtro: procesa todas las lecturas y
calcula dos funciones de ventana, por lo que el tiempo crece con el volumen:
1.423 ms (1k), 13.15 ms (10k), 148.43 ms (100k) y 1\,645.9 ms (1M).

\textbf{Con índices:} el B+ tree \texttt{(id\_sensor, fecha\_hora)} entrega las
filas ya ordenadas y evita el \texttt{Sort}, pero hay que procesar todas las
filas igual. El tiempo queda en 1.395, 13.61, 144.59 y 1\,443.9 ms, con una
mejora de -12.3\,\% en 1M. En 1k el planificador prefiere \texttt{Seq Scan}
(\texttt{usa\_indice = no}), así que la diferencia es ruido. Es el ejemplo de
una consulta donde el índice ayuda poco porque necesita todas las filas.


\subsubsection{Consulta 4}
\resultadoconsulta{4}

\textbf{Sin índices:} por cada alerta grave, el \texttt{JOIN LATERAL} busca las
cinco lecturas anteriores de su sensor; sin índice cada búsqueda recorre toda
la tabla. El tiempo sube a 0.376 ms (1k), 10.70 ms (10k), 1\,114.8 ms (100k) y
92\,572.7 ms (1M).

\textbf{Con índices:} el B+ tree sobre \texttt{(id\_sensor, fecha\_hora)}
permite leer solo las cinco entradas anteriores a cada alerta; el Hash sobre
\texttt{nivel} filtra las alertas graves. El tiempo baja a 0.252, 0.445,
3.552 y 36.33 ms, es decir, alrededor de 2\,548 veces menos con 1M.


\subsubsection{Consulta 5}
\resultadoconsulta{5}

\textbf{Sin índices:} por cada alerta hay que buscar la última lectura normal
anterior a una fecha y comprobar con \texttt{NOT EXISTS} si tiene orden; cada
búsqueda recorre \texttt{LecturaVibracion} completa y el costo crece de forma
cuadrática: 0.458 ms (1k), 74.82 ms (10k), 5\,933.1 ms (100k) y 575\,110.5 ms
(1M, $n=1$).

\textbf{Con índices:} el B+ tree \texttt{(id\_sensor, fecha\_hora)} recorre el
árbol hacia atrás desde la fecha de la alerta y el Hash sobre
\texttt{AlertaOrden(id\_alerta)} resuelve el anti-join. El tiempo baja a
0.060, 0.776, 5.606 y 71.84 ms; con 1M es aproximadamente 8\,005 veces más
rápida, la mayor ganancia del experimento.



\subsection{Análisis y Discusión}


\begin{table}[H]
\centering
\small
\caption{Tiempo promedio con 1\,000\,000 de lecturas y variación al usar índices}
\begin{tabular}{|c|r|r|r|r|}
\hline
\textbf{Consulta} & \textbf{Sin índices (ms)} & \textbf{Con índices (ms)} & \textbf{Variación} & \textbf{Factor} \\ \hline
1 & 197\,558.0 & 55.98 & $\approx-100$\,\% & $\approx3\,529\times$ \\ \hline
2 & 165.97 & 50.54 & $-69.5$\,\% & $\approx3.3\times$ \\ \hline
3 & 1\,645.9 & 1\,443.9 & $-12.3$\,\% & $\approx1.1\times$ \\ \hline
4 & 92\,572.7 & 36.33 & $\approx-100$\,\% & $\approx2\,548\times$ \\ \hline
5 & 575\,110.5 & 71.84 & $\approx-100$\,\% & $\approx8\,005\times$ \\ \hline
\end{tabular}
\end{table}


Las consultas 1, 4 y 5 hacen una búsqueda por cada alerta o lectura anómala.
Sin índices, cada búsqueda recorre la tabla completa y el costo crece como
$N^2$: al pasar de 100k a 1M el tiempo se multiplica por casi 100. Con el B+
tree el acceso es un salto directo y el tiempo crece de forma aproximadamente
lineal, de ahí factores de mejora de miles de veces con 1M.

La consulta 2 filtra por un rango de 48 h. Con pocos datos ese rango abarca casi
toda la tabla y el índice no ayuda; con 1M representa cerca del 4\,\% de las
filas y la ganancia llega a dos tercios. La consulta 3 recorre todas las filas
para calcular las ventanas: el índice solo evita el ordenamiento, por eso la
mejora es pequeña (alrededor de 12\,\% con 1M). Es el caso típico de una
consulta que necesita todas las filas y donde un índice aporta poco.

Con 1k casi todas las consultas están por debajo de unos pocos milisegundos y
las diferencias son ruido. Por eso el análisis se concentra en 100k y 1M.

\paragraph{Costo en escrituras.}
Los índices también cuestan: cada \texttt{INSERT}, \texttt{UPDATE} o
\texttt{DELETE} debe mantener los B+ tree y los Hash. La tabla siguiente lo
cuantifica con 1M de lecturas.

\begin{table}[H]
\centering
\small
\caption{Costo de actualización con 1M (ms; cada operación termina en ROLLBACK)}
\begin{tabular}{|l|r|r|r|}
\hline
\textbf{Operación} & \textbf{Sin índices} & \textbf{Con índices} & \textbf{Variación} \\ \hline
INSERT 1000 lecturas & 24.96 & 16.72 & $-33.0$\,\% \\ \hline
UPDATE rms 1 día & 170.35 & 271.60 & $+59.4$\,\% \\ \hline
UPDATE estado alertas & 7.343 & 7.984 & $+8.7$\,\% \\ \hline
DELETE 1 día $^{\dagger}$ & 958\,585.4 & 5\,956.1 & $-99.4$\,\% \\ \hline
\end{tabular}
\par\footnotesize $^{\dagger}$ Medido una sola vez ($n=1$): el anti-join sin índice tardó cerca de 16 minutos.
\end{table}


El \texttt{INSERT} de 1000 lecturas es más barato con índices porque la
verificación de las llaves foráneas y el trigger encuentran sus filas por Hash
en lugar de recorrer las tablas. El \texttt{UPDATE} de un día de
\texttt{valor\_rms} es más caro con índices (+59.4\,\% de variación), porque hay que
actualizar los B+ tree que incluyen esa columna. Las alertas, al ser una tabla
pequeña, casi no cambian. La diferencia más grande está en el \texttt{DELETE}
con anti-join: sin índice tarda cerca de 16 minutos y con índice unos segundos.

\paragraph{Espacio.}
Con 1M de lecturas, los índices elegidos por teoría ocupan 105.7\,MB, mientras
que la tabla \texttt{LecturaVibracion} ronda los 65\,MB. Es el precio de las
ganancias anteriores: el espacio de los índices puede superar al de la tabla que
los usa.

\paragraph{Escalamiento por encima del millón de registros.}
Las consultas que recorren toda la tabla (como la 3) crecen en proporción al
número de filas; las búsquedas por índice crecen mucho más despacio. Con 10
millones de lecturas, la C5 sin índices pasaría de 9.6 minutos a varias horas,
mientras que con índices seguiría en el orden de segundos. Para volúmenes
mayores conviene particionar \texttt{LecturaVibracion} por fecha, mantener
tablas de resumen y, si el volumen crece mucho más, usar réplicas de lectura.
La arquitectura cliente-servidor con un solo PostgreSQL es suficiente para
algunos millones de registros, pero no para cientos de millones sin
particionamiento.
